\documentclass[10pt,a4paper]{amsart}
\usepackage[margin=19mm]{geometry}
\usepackage{amsmath,amssymb,mathtools}
\usepackage[scr=rsfs,cal=euler]{mathalfa}
\usepackage{xfrac}
\usepackage[unicode,hidelinks,bookmarks=false]{hyperref}
\newcommand{\ie}{i.e.\ }
\newcommand{\cf}{cf.\ }
\newcommand{\resp}{respectively\ }
\newcommand{\Subsection}{\subsection}
\providecommand{\nobreakdash}{\nobreak\hbox{-}}
\setlength{\emergencystretch}{2em}
\multlinegap0pt
\usepackage{bm}
\usepackage{bbm}
\usepackage{dsfont}
\usepackage{xspace}
\usepackage{cancel}
\usepackage{mathtools}
\usepackage{amsthm}
\makeatletter
\@ifundefined{theorem}{\newtheorem{theorem}{Theorem}}{}
\@ifundefined{definition}{\newtheorem{definition}[theorem]{Definition}}{}
\@ifundefined{proposition}{\newtheorem{proposition}[theorem]{Proposition}}{}
\@ifundefined{remark}{\newtheorem{remark}[theorem]{Remark}}{}
\makeatother
\newcommand{\PC}{\operatorname{PC}} % Parking Completions
\newcommand{\PF}{\operatorname{PF}} % Parking Functions
\title[Parking completions are $\mathbf{x}$-parking functions]{Parking completions are $\mathbf{x}$-parking functions}
\author{Kimberly P. Hadaway}
\date{Source: arXiv:2607.16098v1 (CC BY 4.0)}
\begin{document}
\maketitle

\noindent\emph{Source and licence.} Parking completions are $\mathbf{x}$-parking functions. Version arXiv:2607.16098v1 (CC BY 4.0), distributed under the Creative Commons Attribution 4.0 International licence, as recorded in the arXiv licence metadata for this exact version and shown on the abstract page https://arxiv.org/abs/2607.16098v1. Full rights notice: licenses/arxiv-2607.16098-CC-BY-4.0.txt.\par\smallskip

\noindent\textit{Editorial scope.} Counted unit: the Proposition whose source label is \texttt{prop:weakly\_increasing\_sets\_are\_equal} in the article (source0.tex lines 443--446, source label \texttt{prop:weakly\_increasing\_sets\_are\_equal}), delivered together with the article's own complete original proof of it (source0.tex lines 448--479, one \texttt{proof} environment of 32 lines). No mathematics is written, repaired or re-derived: statement and proof are the article's own text line for line. Required context retained uncounted: def:x\_parking\_function (lines 197--199); rem:find\_x\_from\_t (lines 419--431). Every definition, display and label of those lines is retained unchanged. Omitted from this package and not redistributed: 1--196, 200--418, 432--442, 447--447, 480--637. Nothing inside a retained excerpt is omitted or altered: every retained body is delivered exactly as the article's lines are written, so no omission receipt of any kind accompanies this package. Citations: the retained excerpts carry no citation command at all, so no bibliography entry of the article is transcribed and no reference is redistributed. Reference adaptation: none was needed and none was made; every reference inside the delivered unit resolves inside it, so no unresolved reference or citation is produced. Printed slips kept literal: no printed word, symbol, hypothesis or number of the counted statement or of its proof was altered. Layout and class: the article's own class is \texttt{amsart} with packages including fullpage, amssymb, hyperref, amsthm, cleveref, mathtools, listings, xcolor; the wrapper is the lane's pinned \texttt{amsart} editorial wrapper at 10pt on A4, which restates the theorem-like environments the retained text needs and adds the article's own macro definitions that the retained text needs (\texttt{\textbackslash PC}, \texttt{\textbackslash PF}), with extra wrapper packages (bm, bbm, dsfont, xspace, cancel, mathtools). The delivered numbering is the unit's own. Assets: no retained span contains a figure environment, a graphics-inclusion command, a PDF-page inclusion, a table or a TikZ picture; no image file is copied into this package, the package declares no asset dependency and no image file is redistributed with it. Licence: version arXiv:2607.16098v1 of this article is distributed under the Creative Commons Attribution 4.0 International licence, as recorded in the arXiv licence metadata for this exact version and shown on the abstract page https://arxiv.org/abs/2607.16098v1. A full rights notice is delivered at licenses/arxiv-2607.16098-CC-BY-4.0.txt. Characters: every retained excerpt is pure ASCII, so the delivered engine, XeLaTeX with its default Latin Modern text font, reports no missing character.
\begin{definition}\label{def:x_parking_function}
    Let $\mathbb{N}=\{1,2,3,\ldots\}$. For $\mathbf{x} \coloneqq (x_1, x_2, \hdots, x_n) \in \mathbb{N}^n$, an \textbf{$\mathbf{x}$-parking function} is a sequence $\mathbf{a} \in\mathbb{N}^n$ whose increasing rearrangement $\mathbf{a}^\uparrow$ satisfies $\mathbf{a}_i^\uparrow \leq x_1 + x_2 + \cdots + x_i$, for all $i\in[n]$.
\end{definition} 

\begin{remark}\label{rem:find_x_from_t}
    In what follows, the \textbf{content} of a list $\mathbf{a} = (a_1, a_2, \hdots, a_n)$ is the set $\{a_1,a_2,\ldots,a_n\}$, and we denote it $\operatorname{con}(\mathbf{a})$.
    % {1^{\ell_{1}},2^{\ell_{2}},\ldots,n^{\ell_{n}}\}$, where, for all $1 \leq i \leq n$, $\ell_i$ represents the number of occurrences of $i$ in $\alpha$. 
    % We call $\ell_i$ the \textit{multiplicity} of $i$ in $\alpha$.
    Let $\mathbf{c} = (c_1, c_2, \hdots, c_{n-m}) \in [n]^{n-m}$ be a parking completion of length $n$ for $\mathbf{t} = (t_1, t_2, \hdots, t_m)$ as an increasing subset of $[n]$.
    Let $\mathbf{u} \coloneqq (u_1, u_2, \hdots, u_{n-m})$ denote the increasing rearrangement of $[n] \setminus \operatorname{con}(\mathbf{t})$.
    That is, $\operatorname{con}(\mathbf{t}) \cup \operatorname{con}(\mathbf{u}) = [n]$, $\operatorname{con}(\mathbf{t}) \cap \operatorname{con}(\mathbf{u}) = \varnothing$, and $u_i < u_j$ for all $i < j$.    
    Define 
    \[
    \mathbf{x} \coloneqq (u_1, u_2 - u_1, u_3 - u_2, \hdots, u_{n-m} - u_{(n-m)-1}) \in \mathbb{N}^{N}
    \]
    where $N \coloneqq \operatorname{len}(\mathbf{u}) = \operatorname{len}(\mathbf{x}) = n-m$.
\end{remark}

\begin{proposition}\label{prop:weakly_increasing_sets_are_equal}
    The set of weakly increasing parking completions of length $n$ for a fixed $\mathbf{t}$ is equal to the set of weakly increasing $\mathbf{x}$-parking functions of length $n-m$ with $\mathbf{x}$ obtained from $\mathbf{t}$ as described in Remark~\ref{rem:find_x_from_t}.
    Namely, $\PC_{n}^\uparrow(\mathbf{t}) = \PF_{n-m}^\uparrow(\mathbf{x})$ with $\mathbf{t}$ and $\mathbf{x}$ as in Remark~\ref{rem:find_x_from_t}.
\end{proposition}

\begin{proof}
    We proceed via a double containment argument.

    Let $\mathbf{c} \in \PC_n^\uparrow(\mathbf{t})$ for $\mathbf{t} = (t_1, t_2, \hdots, t_m)$.
    We show that $\mathbf{c}^\uparrow \in \PF_{n-m}^\uparrow(\mathbf{x})$.
    Equivalently, we show that $\mathbf{c}^\uparrow$ satisfies $c_i^\uparrow \leq x_1 + x_2 + \cdots + x_i$ for all $i \in [n-m]$.
    Notice $\mathbf{c} = \mathbf{c}^\uparrow$ because we assumed that $\mathbf{c}$ was weakly increasing.
    Now, observe $c_1 \leq x_1 = u_1$ because $u_1$ is the first unoccupied spot on the street in the parking completion setting.
    If $c_1 > u_1 = x_1$, then spot $u_1$ is left empty, and $\mathbf{c} \notin \PC_{n}^\uparrow(\mathbf{t})$ which is a contradiction.
    This is the base case of our inductive argument.
    
    For the induction hypothesis, we assume that if $\mathbf{c} = (c_1, c_2, \hdots, c_k) \in \PC_{m+k}^\uparrow(\mathbf{t})$, then $\mathbf{c} \in \PF_{k}^\uparrow(\mathbf{x})$.
    For the induction step, consider $\mathbf{d} = (d_1, d_2, \hdots, d_k, d_{k+1}) \in \PC_{m+(k+1)}^\uparrow(\mathbf{t})$.
    We show $\mathbf{d} \in \PF_{k+1}^\uparrow(\mathbf{x})$.
    Using the induction hypothesis, $(d_1, d_2, \hdots, d_k) \in \PF_{k}^\uparrow(\mathbf{x})$, and thus, $d_i \leq x_1 + x_2 + \cdots + x_i$ for $i \in [k]$ by definition of $\mathbf{x}$-parking function.
    It remains to check that $d_{k+1} \leq x_1 + x_2 + \cdots + x_{k+1}$.
    Substituting $x_1 = u_1$ and $x_i = u_{i} - u_{i-1}$ for $i \in \{2, 3, \ldots, k+1\}$, we have a telescoping sum that simplifies to show $x_1 + x_2 + \cdots + x_{k+1} = u_{k+1}$.
    We immediately have that $d_{k+1} \leq u_{k+1} = x_1 + x_2 + \cdots + x_{k+1}$ by the inequality characterization of parking completions.
    This shows $\PC_{n}^\uparrow(\mathbf{x}) \subseteq \PF_{n-m}^\uparrow(\mathbf{x})$.

    For the other containment, let $\mathbf{a} = (a_1, a_2, \hdots, a_{n-m}) \in \PF_{n-m}^\uparrow(\mathbf{x})$.
    We show $\mathbf{a} \in \PC_{m+(n-m)}^\uparrow(\mathbf{t})$.
    Equivalently, we show that $\mathbf{a}^\uparrow$ satisfies $a_i^{\uparrow} \leq u_i$ for all $i \in [n]$.
    Notice $\mathbf{a} = \mathbf{a}^\uparrow$ because we assumed that $\mathbf{a}$ was weakly increasing.
    Now, observe $a_1 \leq x_1 = u_1$ by Definition~\ref{def:x_parking_function} and Remark~\ref{rem:find_x_from_t}.
    By Remark~\ref{rem:find_x_from_t}, for all $i \in \{2, 3, \hdots, n-m\}$, we have that $x_i = u_i - u_{i-1}$.
    This implies $a_i \leq x_1 + x_2 + \hdots + x_i = u_i$ for all $i \in \{ 2, 3, \hdots, n-m \}$, as desired.
    We can obtain $\mathbf{u} = (u_1, u_2, \hdots, u_{n-m})$ using the relationships mentioned in the previous two sentences, and we can obtain $\mathbf{t} = (t_1, t_2, \hdots, t_m)$ as the increasing rearrangement of $[n] \setminus \operatorname{con}(\mathbf{u})$, and in this way, $\mathbf{a}$ is a parking completion of length $n$ for $\mathbf{t}$.
    This shows $\PF_{n-m}^\uparrow(\mathbf{x}) \subseteq \PC_{n}^\uparrow(\mathbf{x})$.

    We have shown double containment, and this concludes the proof.
\end{proof}

\end{document}
